\section{Invariant subspaces of the marked rank-three representation}
\label{sec:marked-shape}

% BEGIN SHP1
\paragraph{SHP1. A separating diagonal detects the coordinate subspaces.}
\label{claim:SHP1}
Let $k$ be a field, $H$ a group, and $A(h)$ a family of $m$ by $m$
matrices over $k$. Suppose that two members have difference
$D=\operatorname{diag}(d_0,\ldots,d_{m-1})$ with pairwise distinct
entries. Every common invariant subspace $W\subseteq k^m$ is the
span of a subset of the coordinate vectors. Indeed $D$ preserves
$W$ by subtraction, so each of its spectral projectors
\[
                P_i=\prod_{b\ne i}\frac{D-d_b I}{d_i-d_b}
\]
preserves $W$. For $v\in W$ the vector $P_i v=v_i e_i$ belongs to
$W$. Thus every coordinate vector that occurs with nonzero
coefficient in any member of $W$ belongs to $W$. Conversely those
vectors span $W$, by expanding its elements in the coordinate
basis. In particular, if the family is reducible, a proper nonzero
common invariant subspace supplies a nonempty proper set $E$ of
indices. For every $j\in E$, $i\notin E$ and $h\in H$,
$A(h)_{ij}=0$: the vector $A(h)e_j$ again lies in that subspace.
This argument applies in dimension two as well as in dimension three.
% END SHP1

% BEGIN SHP2
\paragraph{SHP2. The four possible shapes on the sign kernel.}
\label{claim:SHP2}
Let $k$ be a field of characteristic different from two, let
$\eta:G\to\{\pm1\}$ be a character and let $r$ take values in
$\mathrm{GL}_3(k)\rtimes\{\pm1\}$ with this sign. The negative
sign acts on the matrix group by $A\mapsto JA^{-t}J$, where $J$
reverses the basis. Write $A_g=r(g)_{\rm mat}$ and $H=\ker\eta$.
Suppose $\eta(c)=-1$, $g_1\in H$, $A_c=I$, and the $(0,1)$ and
$(1,0)$ entries of $Z=A_{g_1}-A_{g_1^*}$ vanish, where
$g^*=cg^{-1}c$. Assume $z=Z_{00}\ne0$. Under these hypotheses
at least one of the following holds on $H$: the three-dimensional
matrices are irreducible (Type I), all are upper triangular, all
are lower triangular, or they preserve the middle coordinate and
the outer two-coordinate plane separately with an irreducible action
on the latter (Type II).

For $g\in H$ the group law gives $g^*\in H$ and
$A_{g^*}=JA_g^tJ$. Thus $Z_{ij}=(A_{g_1})_{ij}-
(A_{g_1})_{2-j,2-i}$. Positions fixed by the reflection
$(i,j)\mapsto(2-j,2-i)$ have zero difference; the others have
opposite reflected entries. The two marked zeros, with their
reflections, give $Z=\operatorname{diag}(z,0,-z)$, which
has distinct diagonal entries by the assumptions on $z$ and $k$.
Moreover a common zero at $(i,j)$ on $H$ forces a common zero
at $(2-j,2-i)$, by applying it to $g^*$ for every $g\in H$.

If the three-dimensional family is reducible, apply SHP1 to this
marked difference. For each possible invariant coordinate subset
$E$, the following are the zeros provided by invariance before
using reflection. A pair $ij$ denotes row $i$, column $j$.
\[
\begin{array}{c|c|l}
 E & \text{forced zeros} & \text{shape after reflection}\\ \hline
 \{0\}   &10,20& \text{upper}\\
 \{0,1\} &20,21& \text{upper}\\
 \{2\}   &02,12& \text{lower}\\
 \{1,2\} &01,02& \text{lower}\\
 \{1\}   &01,21& \text{middle line and outer plane}\\
 \{0,2\} &10,12& \text{middle line and outer plane}
\end{array}
\]
Indeed reflection pairs $10$ with $21$ and $01$ with $12$. In
either of the last two rows all four adjacent entries vanish.
The action on the outer plane has marked difference
$\operatorname{diag}(z,-z)$. If it is irreducible, the precise
Type II condition follows. If it is reducible, SHP1 in dimension
two makes one of the lines $k e_0,k e_2$ invariant. In the first
case the entry $20$ also vanishes throughout and the full family
is upper triangular; in the second the entry $02$ vanishes and the
full family is lower triangular. This proves the claimed disjunction
over the original field without a change of basis.

For the points of VSCM8 the hypotheses follow in the same basis:
integral marking and its two zeros persist under torus conjugation,
and VSCM6 gives $|z|=1$ from the marked unit at the residual point.
For a two-sided point VSCM9 excludes both triangular alternatives.
The remaining distinction between Type I and Type II requires the
arithmetic exclusion of the irreducible outer block.
% END SHP2
